\documentclass[11pt, a4paper]{article}

\usepackage{ctex}
\usepackage{geometry}
\usepackage{titlesec}
\usepackage{enumitem}
\usepackage{setspace}
\usepackage{hyperref}


\title{Script 大纲}
\author{黄浩宇}
\date{\today}

\begin{document}

        \maketitle
        \newpage
        \tableofcontents

        \newpage

        \section{【剧本大纲】}

                \subsection{主题立意}
                        \begin{itemize}[leftmargin=2em, itemsep=0.2em]
                                \item \textbf{核心主题}：讽刺教育领域中的形式主义、统一口径与虚假繁荣。
                                \item \textbf{表现手法}：以机器人般的动作、整齐划一的掌声与笑声、新闻宣传等元素，制造荒诞感与压迫感。
                                \item \textbf{情感基调}：冷峻、荒诞、讽刺，带有黑色幽默色彩。
                        \end{itemize}

                \subsection{人物设置}
                        \begin{itemize}[leftmargin=2em, itemsep=0.2em]
                                \item \textbf{校长}：学校最高决策者，形式主义的推行者。
                                \item \textbf{开明老师}：唯一提出异议的人，代表被压制的理性声音。
                                \item \textbf{教育部部长}：上级视察领导，只接受符合预期的回答。
                                \item \textbf{秘书}：部长的随行人员，负责“处理”不合口径的学生。
                                \item \textbf{学生A、B、C}：面对视察的不同反应者。A、B说真话被拉走，C说官话被表扬。
                                \item \textbf{群众领导、群众老师}：随行听会人员，像机器一样附和。
                                \item \textbf{群众学生}：学生群体，像机器人一样学习与回答。
                        \end{itemize}

                \subsection{结构脉络}
                        \begin{enumerate}[leftmargin=2em, itemsep=0.3em]
                                \item \textbf{开场}：鲨鱼血腥嘴巴与尖牙，伴随凄厉惨叫，奠定压抑荒诞的基调。
                                \item \textbf{第1幕：会议室}：校长提出“提高综合素养”“周六上课、开设兴趣班”。全场像机器一样整齐鼓掌、整齐发笑；开明老师提出异议，被校长用荒谬逻辑反问并遭全场嘲笑。校长宣布散会，所有人像机器人一样走出会场。
                                \item \textbf{第2幕：教室}：星期六，学生像机器人一样参加“电路设计兴趣小组”。报时器一响，学生整齐通电、亮灯，并异口同声承认“素质提高了”；随后整齐离开。黑板上字变为“兴趣小组结束”，相机快门响起，学校网站发布新闻，宣称该校成为全市领先高中。
                                \item \textbf{第3幕：学校走廊 / 校园内}：教育部部长视察，随机采访学生。学生A、B说“很无聊，听不懂”，被秘书拉走；学生C用标准官话称赞活动“非常有意义”“受益匪浅”，被部长表扬。随后，本市教育部发布新闻，宣称全市学生素质普遍提高、认可度极高，已成为全省教育先进市。
                        \end{enumerate}

                \subsection{关键冲突}
                        \begin{itemize}[leftmargin=2em, itemsep=0.2em]
                                \item \textbf{校长 vs 开明老师}：形式主义要求与理性异议之间的冲突。
                                \item \textbf{部长 vs 学生A、B}：真实感受与官方期待之间的冲突。
                                \item \textbf{学生C vs 学生A、B}：迎合口径与说真话之间的对比。
                                \item \textbf{新闻宣传 vs 现实情况}：表面繁荣与真实体验之间的反差。
                        \end{itemize}

                \subsection{视听设计}
                        \begin{itemize}[leftmargin=2em, itemsep=0.2em]
                                \item \textbf{动作}：机器人般的步伐、整齐划一的鼓掌与笑声。
                                \item \textbf{音效}：鲨鱼惨叫、报时器“滴”声、相机快门声。
                                \item \textbf{画面}：电子钟显示星期六、黑板字自动改名、新闻内容弹出。
                                \item \textbf{节奏}：从会议室的压抑，到教室的机械，再到视察的荒诞，层层递进。
                        \end{itemize}

                \subsection{结尾处理}
                        \begin{itemize}[leftmargin=2em, itemsep=0.2em]
                                \item 以本市教育部发布新闻收束，形成“现实被宣传覆盖”的闭环。
                                \item 不直接给出批判结论，而是让荒诞感停留在新闻口号与机器人动作的对照中。
                        \end{itemize}

\end{document}